\documentclass[a4paper,11pt]{article}
\usepackage{midterm-style}
\examinehead{企业实训课程 · 期中考试}{D 卷}
\begin{document}

\examcover{企业实训课程}{2026—2027 学年第一学期期中考试（D 卷）}{计算机组成原理、操作系统、计算机网络、数据库概论、信息系统导论（第 1—6 章）}{150 分}{180 分钟}{本卷为开卷考试，允许携带教材、笔记和纸质资料，可使用手机或电脑检索公开资料；不得相互讨论、传递答案，不得使用手机通讯或社交软件。所有答案写在答题纸上，计算题须写出主要步骤。}

% ============================================================
\parttitle{一、计算机组成原理（25 分）}
% ============================================================

\qtypename{（一）单项选择题}{5 题，每题 2 分，共 10 分}

\q 冯·诺依曼瓶颈主要是指（\quad）
\choices{A. CPU 与主存之间的速度差异}{B. CPU 与 Cache 之间的容量差异}{C. 外设与主存之间的接口差异}{D. 总线与寄存器之间的宽度差异}

\q 8 位补码 \code{01111111} 加 \code{00000001}，结果和溢出标志 OF 分别为（\quad）
\choices{A. \code{10000000}，OF=1}{B. \code{10000000}，OF=0}{C. \code{00000000}，OF=1}{D. \code{11111111}，OF=0}

\q 下列关于 DRAM 的叙述中，正确的是（\quad）
\choices{A. 不需要刷新，速度快}{B. 需要定期刷新，集成度高}{C. 用触发器存储信息}{D. 断电后数据不会丢失}

\q 某总线宽度为 32 位，时钟频率为 100 MHz，每个时钟周期传送一次数据，则总线带宽约为（\quad）
\choicesline{A. 100 MB/s}{B. 200 MB/s}{C. 400 MB/s}{D. 800 MB/s}

\q CPU 响应可屏蔽中断的条件不包括（\quad）
\choices{A. 中断源有中断请求}{B. CPU 允许中断}{C. 当前指令执行结束}{D. 当前进程必须处于阻塞态}

\needspace{30\baselineskip}
\qtypename{（二）基础计算题}{2 题，共 8 分}

\needspace{18\baselineskip}
\q （8 分）完成下列计算，写出必要过程：
\begin{enumerate}[label=（\arabic*）, leftmargin=2.4em, itemsep=1pt]
  \item （4 分）将二进制数 $1110.011_2$ 分别转换为十进制数和十六进制数。
  \item （4 分）将十进制数 $25.625$ 转换为 IEEE 754 单精度（32 位）浮点数，写出其二进制格式和十六进制表示。
\end{enumerate}
\anslines{5.5cm}

\qtypename{（三）简答题}{1 题，共 7 分}

\needspace{16\baselineskip}
\q （7 分）简述 Cache 的工作原理，写出平均访问时间的计算公式，并比较写直达（write-through）与写回（write-back）两种写策略的优缺点。
\anslines{6cm}

% ============================================================
\parttitle{二、操作系统（20 分）}
% ============================================================

\qtypename{（一）单项选择题}{4 题，每题 2 分，共 8 分}

\q 下列进程通信方式中，通常速度最快的是（\quad）
\choicesline{A. 管道}{B. 消息队列}{C. 共享内存}{D. 套接字}

\q 高响应比优先调度算法的主要优点是（\quad）
\choices{A. 只照顾短作业}{B. 只照顾长作业}{C. 兼顾长短作业，减少饥饿}{D. 不需要计算响应比}

\q 最佳适应（Best Fit）分区分配算法的主要缺点是（\quad）
\choices{A. 产生大量难以利用的小碎片}{B. 总是分配最大分区}{C. 查找速度最快}{D. 不会产生外部碎片}

\q 文件系统中，索引结点（inode）主要存放（\quad）
\choices{A. 文件的内容数据}{B. 文件的元数据和数据块地址}{C. 磁盘分区表}{D. CPU 寄存器内容}

\needspace{30\baselineskip}
\qtypename{（二）简答题}{1 题，共 12 分}

\needspace{28\baselineskip}
\q 回答下列问题：
\begin{enumerate}[label=（\arabic*）, leftmargin=2.4em, itemsep=2pt]
  \item （6 分）简述虚拟存储器的基本思想，并说明请求分页系统中缺页中断的处理过程。
  \item （6 分）某系统有 5 个进程 $P_0\sim P_4$，资源类型为 A、B、C。当前 Available $=(3,3,2)$，各进程的 Max 和 Allocation 如下表：
  \begin{center}\small
  \begin{tabular}{ccccc}
  \toprule
  进程 & Max(A,B,C) & Allocation(A,B,C) & Need(A,B,C) \\
  \midrule
  $P_0$ & (7,5,3) & (0,1,0) & \\
  $P_1$ & (3,2,2) & (2,0,0) & \\
  $P_2$ & (9,0,2) & (3,0,2) & \\
  $P_3$ & (2,2,2) & (2,1,1) & \\
  $P_4$ & (4,3,3) & (0,0,2) & \\
  \bottomrule
  \end{tabular}
  \end{center}
  请计算各进程的 Need，并给出一个安全序列。
\end{enumerate}
\anslines{9cm}

% ============================================================
\parttitle{三、计算机网络（35 分）}
% ============================================================

\qtypename{（一）单项选择题}{5 题，每题 2 分，共 10 分}

\q 某链路带宽为 100 Mb/s，往返时延 RTT 为 10 ms，则时延带宽积约为（\quad）
\choicesline{A. 1 Mb}{B. 125 KB}{C. 1.25 MB}{D. 12.5 KB}

\q 以太网交换机通过自学习建立转发表时，主要根据（\quad）
\choicesline{A. 目的 MAC 地址}{B. 源 MAC 地址}{C. 源 IP 地址}{D. 目的端口号}

\q 下列 IP 地址中属于私有地址的是（\quad）
\choicesline{A. 8.8.8.8}{B. 172.16.0.1}{C. 192.169.1.1}{D. 11.0.0.1}

\q 下列关于 UDP 的叙述中，正确的是（\quad）
\choices{A. 面向连接、可靠传输}{B. 无连接、不可靠，支持广播和多播}{C. 提供拥塞控制}{D. 使用三次握手建立连接}

\q FTP 工作时通常使用的 TCP 连接数是（\quad）
\choicesline{A. 1 个}{B. 2 个}{C. 3 个}{D. 4 个}

\needspace{30\baselineskip}
\qtypename{（二）综合计算与应用题}{3 题，共 25 分}

\needspace{24\baselineskip}
\q （8 分）某单位分配到地址块 192.168.1.0/24，需要划分给 6 个部门，每个部门至少 25 台主机。请回答：
\begin{enumerate}[label=（\arabic*）, leftmargin=2.4em, itemsep=1pt]
  \item 应选择多长的子网前缀？为什么不能用 /28？
  \item 写出前 6 个子网的地址/掩码、可用主机地址范围和广播地址。
  \item 剩余地址空间还能划分出多少个同样大小的子网？
\end{enumerate}
\anslines{8cm}

\needspace{24\baselineskip}
\q （8 分）某 TCP 连接的拥塞窗口 \code{cwnd} 随 RTT 变化如下：
\[
\begin{array}{c|ccccccccccccccccc}
\text{RTT} & 1&2&3&4&5&6&7&8&9&10&11&12&13&14&15&16&17\\
\hline
\text{cwnd/MSS} & 1&2&4&8&9&10&11&12&13&14&15&16&1&2&4&8&9
\end{array}
\]
请回答：
\begin{enumerate}[label=（\arabic*）, leftmargin=2.4em, itemsep=1pt]
  \item 初始慢开始门限 \code{ssthresh} 约为多少？说明依据。
  \item 第 13 个 RTT 可能发生了什么事件？此时 \code{ssthresh} 和 \code{cwnd} 如何变化？
  \item 第 17 个 RTT 时若收到 3 个重复 ACK，按快重传/快恢复，\code{ssthresh} 和 \code{cwnd} 变为多少？收到新 ACK 后 \code{cwnd} 又为多少？
  \item 指出哪些区间处于慢开始、哪些区间处于拥塞避免。
\end{enumerate}
\anslines{8cm}

\needspace{24\baselineskip}
\q （9 分）某小型企业从运营商获得公网地址 202.1.1.2/30，运营商网关为 202.1.1.1；企业内部使用 192.168.1.0/24，内网 Web 服务器为 192.168.1.100:80。请说明：
\begin{enumerate}[label=（\arabic*）, leftmargin=2.4em, itemsep=1pt]
  \item 内网 PC 的 IP 地址、子网掩码、默认网关和 DNS 应如何配置？
  \item 若要让内网 PC 访问 Internet，路由器需要配置什么功能？
  \item 若要让外网访问内网 Web 服务器，应如何配置端口映射？写出外网用户访问的地址。
  \item 简述外网请求到达 Web 服务器过程中源/目的地址的转换过程。
\end{enumerate}
\anslines{8cm}

% ============================================================
\parttitle{四、数据库概论（35 分）}
% ============================================================

\qtypename{（一）单项选择题}{5 题，每题 2 分，共 10 分}

\q 数据库三级模式结构中，描述数据库中全体数据全局逻辑结构的是（\quad）
\choicesline{A. 外模式}{B. 模式}{C. 内模式}{D. 存储模式}

\q 设关系 $R(A,B)$ 和 $S(B,C)$，则 $R$ 与 $S$ 自然连接后结果关系的属性个数为（\quad）
\choicesline{A. 2}{B. 3}{C. 4}{D. 5}

\q 设某列有 10 行，其中 3 行为空值，则 \code{COUNT(列名)} 的结果是（\quad）
\choicesline{A. 10}{B. 7}{C. 3}{D. 0}

\q 下列关于 BCNF 的叙述中，正确的是（\quad）
\choices{A. BCNF 一定属于 3NF}{B. 3NF 一定属于 BCNF}{C. BCNF 允许非主属性对候选码的传递依赖}{D. BCNF 不允许任何函数依赖}

\q 能避免脏读、但不能避免不可重复读的事务隔离级别是（\quad）
\choicesline{A. 读未提交}{B. 读已提交}{C. 可重复读}{D. 串行化}

\needspace{30\baselineskip}
\qtypename{（二）综合题}{3 题，共 25 分}

\needspace{30\baselineskip}
\q （10 分）设学生选课数据库中有三个基本表（加下划线的属性为主码）：
\begin{center}\small
\begin{tabular}{ll}
\toprule
关系名 & 属性 \\
\midrule
Student & \underline{Sno}, Sname, Sdept, Sage \\
Course & \underline{Cno}, Cname, Ccredit \\
SC & \underline{Sno}, \underline{Cno}, Grade \\
\bottomrule
\end{tabular}
\end{center}
请用标准 SQL 完成下列各题：
\begin{enumerate}[label=（\arabic*）, leftmargin=2.4em, itemsep=1pt]
  \item （2.5 分）查询选修了全部课程的学生学号和姓名。
  \item （2.5 分）查询每门课程的不及格人数（成绩小于 60 分）。
  \item （2.5 分）查询平均成绩在 80 分以上的课程号和平均成绩。
  \item （2.5 分）创建 \code{Student} 表，要求 \code{Sno} 为主码，\code{Sname} 非空，\code{Sage} 在 15～60 之间，\code{Sdept} 默认值为“未分配”。
\end{enumerate}
\anslines{9cm}

\needspace{24\baselineskip}
\q （8 分）某电商订单系统中有关系模式
\[
R(\text{订单号},\text{客户号},\text{客户名},\text{商品号},\text{商品名},\text{单价},\text{数量}),
\]
其函数依赖集为
\[
F=\{\text{订单号}\to\text{客户号},\ \text{客户号}\to\text{客户名},\ \text{商品号}\to\text{商品名},\ \text{商品号}\to\text{单价},\ (\text{订单号},\text{商品号})\to\text{数量}\}.
\]
请回答：
\begin{enumerate}[label=（\arabic*）, leftmargin=2.4em, itemsep=1pt]
  \item 写出 $R$ 的全部候选码。
  \item $R$ 最高属于第几范式？说明理由。
  \item 将 $R$ 分解为 3NF，要求分解保持函数依赖并具有无损连接性，并简要说明理由。
\end{enumerate}
\anslines{7cm}

\needspace{24\baselineskip}
\q （7 分）说明事务的四种隔离级别（读未提交、读已提交、可重复读、串行化）分别能避免哪些并发异常（脏读、不可重复读、幻读），并说明隔离级别与并发性能之间的关系。
\anslines{7cm}

% ============================================================
\parttitle{五、信息系统导论（第 1—6 章）（35 分）}
% ============================================================

\qtypename{（一）名词解释}{3 题，每题 5 分，共 15 分；要求写出概念内涵，必要时举例}

\q 数据
\anslines{3cm}

\q 信息系统
\anslines{3cm}

\q 软件工程
\anslines{3cm}

\needspace{30\baselineskip}
\qtypename{（二）简答题}{1 题，共 8 分}

\needspace{16\baselineskip}
\q （8 分）简述大数据的 3V/4V 特征，并说明 Hadoop 的主要组成部分及作用。
\anslines{5.5cm}

\needspace{36\baselineskip}
\qtypename{（三）材料分析题}{1 题，共 12 分}

\q （12 分）阅读下列材料，回答问题。
\begin{quote}\small
某城市建设智能交通系统：路口部署摄像头、地磁和雷达传感器采集车流；数据经 5G 或光纤传至交通云平台；云平台保存车辆、违法和设备数据，运行流量分析模型，动态调整信号灯配时，并通过手机 APP 发布路况；部分路口部署边缘计算节点，断网时仍能本地控制信号灯。
\end{quote}
\begin{enumerate}[label=（\arabic*）, leftmargin=2.4em, itemsep=1pt]
  \item （4 分）按“感知层—网络层—平台层—应用层”四层结构分析该系统，说明各层的主要功能。
  \item （4 分）用数据、信息、知识、智慧的层次关系说明车流数据如何转化为信号灯配时决策，并指出数据库、计算机网络、云计算和边缘计算的作用。
  \item （4 分）从系统可靠性和数据安全角度，各提出一条改进建议并说明理由。
\end{enumerate}
\anslines{12cm}

\vfill
\begin{center}
\textbf{—— 试卷正文结束，附页见后 ——}
\end{center}

\answerappendix{附页（一）}{如答题纸空间不足，可在此页继续作答，并注明题号。}
\answerappendix{附页（二）}{如答题纸空间不足，可在此页继续作答，并注明题号。}

\end{document}
